\section{Topological Groups}
\begin{defi}
    Given an abelian group $(G,+)$ and a topology on $G$ we say that $G$ is a \emph{topological abelian group} if $+\colon G\times G\to G$ and $-\colon G\to G$ are continuous.
\end{defi}
Note that if $\mathcal B_0$ is a local basis at 0 then $T_g[\mathcal B_0]$ is a local basis at $g\in G$, where $T_g\colon G\to G$ is given by $T_g(h)=g+h$ which is a homeomorphism. That is, the topology of $G$ is determined by $\mathcal B_0$.

\begin{thm}\label{thm1}
Let $H$ be the intersection of all neighborhoods of $0$. Then $H$ is a subgroup of $G$ and can be described as the closure of $\{0\}$. Furthermore, $G$ is Hausdorff if and only if $H=0$.    
\end{thm}
\begin{proof}
    By construction $0\in H$. Consider the continuous function $f\colon (g,h)\in G\times G\mapsto g-h\in G$. Given $g,h\in H$ and $U$ a neighborhood of $0$. By continuity, $f(0,0)=0$ and the product topology, there exists $V$ and $W$ neighborhoods of $0$ such that that $V\times W\subseteq f^{-1}[U]$. Therefore $g-h=f(g,h)\in f[V\times W]\subseteq f[f^{-1}[U]]\subseteq U$. Since $U$ was arbitrary and by the subgroup criterion we have $H\leq G$. On the other hand,
    \begin{align*}
             g\in\overline{\{0\}} &\iff 0\in U \quad \forall U\text{ neighborhood of }g \\
             &\iff 0\in g+U \quad \forall U\text{ neighborhood of }0 \\
             &\iff g\in H.
    \end{align*}
    So $H=\overline{\{0\}}$. If $G$ is Hausdorff all finite sets are closed and by definition of closure $H=0$. Recall that $G$ is Hausdorff iff $\Delta_G$ is closed. Thus, if $H=0$, since $\Delta_G=f^{-1}[\{0\}]$, $G$ is Hausdorff.
\end{proof}

\subsection{Topology generated by a decreasing chain}

Let $(G,+)$ denote an abelian group. Consider a decreasing chain of subgroups of $G$,
\[
G=G_0\supseteq G_1\supseteq G_2\supseteq \cdots \supseteq G_n\supseteq \cdots
\]
such chain will be called \emph{fundamental system of neighborhoods of 0}. The reason behind this name is shown in the following 
\begin{thm}
    $\mathcal B:= \{T_g\left[G_n\right]: n\geq 0, \, g\in G\}\subseteq G$ is a basis of a topology on $G$ such that $G$ becomes a topological abelian group.
\end{thm}
\begin{proof}
    Let us first see that $\mathcal B$ is a basis of a topology on $G$. Indeed: $G=\bigcup\mathcal B$ and $k+G_n\subseteq T_g[G_n]\cap T_h[G_m]$ if $m\leq n$ and $k\in T_g[G_n]\cap T_h[G_m]$. To see that it is a topological abelian group it suffices to verify that $+$ and $-$ are continuous. Since $T_h(G_n)=T_g(G_n)$ if and only if $g-h\in G_n$, $T_g[G_n]+T_h[G_n]=T_{g+h}[G_n]$ and $-T_g[G_n]=T_{-g}[G_n]$ we are done.
\end{proof}

Note that the topological abelian subgroups $G/G_n$ with respect the quotient topology are discrete ---in fact $G/H$ is discrete iff $H$ is open--- and can be viewed as the topology given by the sequence $(G_m/G_n)_{m\leq n}$; i.e., the image of the fundamental system of neighborhoods of $0$ of $G$ via the canonical projection.
% From now on we will denote $\mathcal B_g=\{T_g[G_n] : n\geq 0\}$ a local basis at $g\in G$.

\subsection{Cauchy completion}

Recall that if $(X,\tau)$ is a topological space and $(x_n)_{n\geq 0}\in X^\mathbb N$ a sequence, we say that $x_n$ converges to $x\in X$ ---written $x_n\to x$--- if for every neighborhood $U$ of $x$ there exists $n_U\in\mathbb N$ such that $x_n\in U$ for $n\geq n_U$.
\begin{defi}
    Let $G$ be a topological abelian group. A sequence $(g_n)_{n\geq 0}\in G^\mathbb N$ is a \emph{Cauchy sequence} if for every neighborhood $U$ of 0 there exists $n_U\in\mathbb N$ such that $x_n-x_m\in U$ for every $n,m\geq n_U$. We denote by $G^\mathbb N_{\mathcal C}$ the set of all Cauchy sequences of $G$.
\end{defi}

Define an operation in $G^\mathbb N_\mathcal C$ component wise and we get an abelian group structure since the sum of two Cauchy sequences is a Cauchy sequence. 
\begin{defi}
    Two Cauchy sequences $(g_n)_{n\geq 0}$, $(h_n)_{n\geq 0}\in G^\mathbb N_\mathcal C$ are \emph{equivalent} if $g_n-h_n\to 0$ in $G$.
\end{defi}

It is clear that equivalence class $[(0)_{n\geq 0}]$ is a subgroup of $G^\mathbb N_\mathcal C$. We denote by $\hat G:=G^\mathbb N_\mathcal C/[(0)_{n\geq 0}]$ the quotient abelian group. For each $g\in G$ we can define a constant sequence $(g)_{n\geq 0}\in G^\mathbb N_\mathcal C$ and take it's equivalence class given rise to a homomorphism
\[
\begin{matrix}
\phi\colon & G&\to & \hat G \\
& g &\mapsto & [(g)_{n\geq 0}].
\end{matrix}
\]
\begin{thm}\label{Hausdorff}
    $\phi$ is one-to-one if and only if $G$ is Hausdorff.
\end{thm}
\begin{proof}
        Since $G\hookrightarrow G^\mathbb N_\mathcal C$ we only need to check the homomorphism $p\colon G^\mathbb N_\mathcal C\twoheadrightarrow \hat G$. It is clear that $\ker p=[(0)_{n\geq 0}]$. Hence, by theorem \ref{thm1},
        \begin{align*}
            \phi \text{ is one-to-one}\iff[(0)_{n\geq 0}]=0 &\iff \overline{\{0\}}=\{0\} \iff G \text{ is Hausdorff}.
        \end{align*}
\end{proof}

If the topology of $G$ is given by a fundamental system of neighborhoods of $0$, $(G_n)_{n\geq 0}$ then we can endow $\hat G$ with a topology giving by the decreasing chain
\[
\hat G=\hat G_0\supseteq \hat G_1\supseteq \hat G_2\supseteq \cdots \supseteq \hat G_n\supseteq \cdots
\]
where $\hat G_n:=\phi[G_n]$; i.e., the equivalence class of sequences that eventually fall inside $G_n$. It is easy to see that $\hat G$ is a topological abelian group and $\phi$ is continuous with respect this choice of topology. By theorem \ref{thm1} and the definition of $\hat G_n$, since
\[
\bigcap_{U\text{ neighborhood of 0}} U=\bigcap_{n\geq 0} \hat G_n=0 
\]
we deduce that $\hat G$ is Hausdorff.

To finish this subsection we are going to see the functoriality properties of $\,\hat{}\,$. Let $f\colon G\to H$ be a continuous homomorphism where $G$ and $H$ are topological abelian groups. Given a Cauchy sequence $(g_n)_{n\geq 0}$, consider the sequence $(f(g_n))_{n\geq 0}$. Let $U$ be any neighborhood of $0$ of $H$. Since $f$ is continuous and $f(0)=0$, $f^{-1}[U]$ is a neighborhood of $0$ in $G$. Thus there exists an integer $n_U$ such that $g_n-g_m\in f^{-1}[U]$ for all $n,m\geq n_U$ so 
\[
f(g_n)-f(g_m)=f(g_n-g_m)\in f[f^{-1}[U]]\subseteq U
\]
for all $n,m\geq n_U$. Hence $f$ induces a homomorphism $\hat f\colon \hat G\to\hat H$. It is easy to check that $\widehat{f\circ g}=\hat f\circ \hat g$ and the following diagram commutes
\[
\begin{tikzcd}
    G \ar[r, "f"] \ar[d, "\phi_G"']& H \ar[d, "\phi_H"]\\
    \hat G \ar[r, "\hat f"] & \hat H.
\end{tikzcd}
\]
 
\subsection{Inverse limit as an algebraic completion}
\begin{defi}
    An \emph{inverse limit} is a limit indexed by the category $\omega^{\text{op}}$; i.e., the terminal cone of a diagram
    \begin{equation}\label{diag1}
    \begin{tikzcd}
     \cdots\ar[r, "\theta_{4}"] & G_3 \ar[r, "\theta_{3}"] & G_2 \ar[r, "\theta_{2}"] & G_1 \ar[r, "\theta_{1}"] & G_0
    \end{tikzcd}
    \end{equation}
    which is denoted as $\varprojlim G_n$:
    \[
    \begin{tikzcd}
    \varprojlim G_n \ar[dd, "\cdots"] \ar[ddr, "\pi_3" description] \ar[ddrr, "\pi_2" description] \ar[ddrrr, "\pi_1" description] \ar[ddrrrr, "\pi_0" description]\\
    {}\\
     \cdots\ar[r, "\theta_{4}"'] & G_3 \ar[r, "\theta_{3}"'] & G_2 \ar[r, "\theta_{2}"'] & G_1 \ar[r, "\theta_{1}"'] & G_0
    \end{tikzcd}
    \]
\end{defi}

We will only consider diagrams $F\colon \omega^{\text{op}}\to \mathbf{Ab}$ called \emph{inverse systems}, which consist of a sequence $(G_n)_{n\geq 0}$ of abelian groups and homomorphisms $\theta_{n+1}\colon G_{n+1}\to G_n$ for every $n$. If each $\theta_n$ is surjective the inverse system is called \emph{surjective}. Note that, in this case, $\pi_n$ ---defined above--- are also surjective.
\begin{thm}\label{varprojalg}
    We have 
    \[
    \varprojlim G_n=\left\{(g_n)_{n\geq 0}\in \prod_{n\geq 0}G_n : \theta_{n+1}(g_{n+1})=g_n \,\,\,\forall n\right\}.
    \]
\end{thm}
\begin{proof}
    Let us denote $\widetilde G$ the set defined on the right hand side. If $A$ is a cone of the diagram \ref{diag1}; i.e.,
    \[
    \begin{tikzcd}
    A \ar[dd, "\cdots"] \ar[ddr, "\mu_3" description] \ar[ddrr, "\mu_2" description] \ar[ddrrr, "\mu_1" description] \ar[ddrrrr, "\mu_0" description]\\
    {}\\
     \cdots\ar[r, "\theta_{4}"'] & G_3 \ar[r, "\theta_{3}"'] & G_2 \ar[r, "\theta_{2}"'] & G_1 \ar[r, "\theta_{1}"'] & G_0
    \end{tikzcd}
    \]
    Then the homomorphism $f\colon a\in A \mapsto (\mu_n(a))_{n\geq 0}\in \widetilde G$ is unique by extensionality and makes $\widetilde G$ a terminal cone of \eqref{diag1}.
\end{proof}
\begin{defi}
     A sequence $(g_n)_{n\geq 0}$ is \emph{coherent} if $(g_n)_{n\geq 0}\in\varprojlim G_n$.
\end{defi}

Let's study the case where $G$ is a topological abelian group given by the sequence $(G_n)_{n\geq 0}$ and consider the inverse limit of the inverse system\footnote{Note that $G/G_n\cong (G/G_{n+1})/(G_n/G_{n+1})$.}
\[
\begin{tikzcd}
       \cdots \ar[r, "p_{n+2}", twoheadrightarrow] &G/G_{n+1} \ar[r, "p_{n+1}", twoheadrightarrow] & G/G_n \ar[r, "p_{n}", twoheadrightarrow] & \cdots \ar[r, "p_{3}", twoheadrightarrow] & G/G_2 \ar[r, "p_{2}", twoheadrightarrow] & G/G_1 \ar[r, "p_1", twoheadrightarrow] & G
    \end{tikzcd}
\]

We endow $\varprojlim G/G_n$ with the initial topology generated by $\{\pi_m\colon \varprojlim G/G_n\to G/G_m\}_{m\geq 0}$ ---the topology generated by the subbase $\{\pi_n^{-1}[U] : U\text{ is open in }G/G_n \}_{n\geq 0}$--- so that it becomes a topological abelian group\footnote{It is not a "straight-forward-proof" but it will be done using an homeomorphic isomorphism with the topological abelian group $\hat G$.}. 

\subsection{Cauchy and inverse limit comparison}

We would like to compare $\varprojlim G/G_n$ with $\hat G$.\todo{Maybe write it with theorems instead of text} It is clear that given $(g_n)_{n\geq 0}\in G^\mathbb N_\mathcal C$ then its image under $p\colon G\to G/G_m$ is ultimately constant; i.e., given the neighborhood $G_m$ since $(g_n)_{n\geq 0}$ is Cauchy, there exists an $n_m\in\mathbb N$ such that $g_k-g_l\in G_m$ for $k,l\geq n_m$ that is $g_k+G_m=g_l+G_m$. Let's denote this constant by $\gamma^{(g_n)_{n\geq 0}}_m\in G/G_m$---or $\gamma_m$ when $(g_n)_{n\geq 0}$ is clear from context---. Since $G_{m+1}\leq G_m$ the sequence $(\gamma_m)_{m\geq 0}$ is coherent. By the definition of equivalence of Cauchy sequences the choice of $\gamma_m$ is well defined in $\hat G$. Hence we have a map
\[
\begin{matrix}
\varphi\colon &\hat G&\longhookrightarrow &\varprojlim G/G_n \\
&&&\\
&[(g_n)_{n\geq 0}] &\longmapsto & (\gamma^{(g_n)_{n\geq 0}}_m)_{m\geq 0}.
\end{matrix}
\]
Conversely, given $(\gamma_n)_{n\geq 0}\in\varprojlim G/G_n$ we can choose a $g_n\in G$ for each $n\geq 0$ such that $\gamma_n=g_n+G_n$. By the definition of coherence and theorem \ref{varprojalg} we have $g_{n+1}+G_n=g_n+G_n$ therefore, $(g_n)_{n\geq 0}\in G^\mathbb N_\mathcal C$. Thus, $\varphi([(g_n)_{n\geq 0}])=(\gamma_n)_{n\geq 0}$, concluding that $\varphi$ is an isomorphism. It is easy to see that for any map $\pi_n$ the composition $\pi_n\circ\varphi$ is constant hence continuous, by the universal property of the initial topology on $\varprojlim G/G_n$, implies that $\varphi$ is continuous. On the other hand, the image of $\hat G_m$ under $\varphi$ are the sets 
\[
\{(\gamma_n)_{n\geq 0}\in\varprojlim G/G_n : \gamma_j=0 \quad \forall j\leq m\}=\bigcap_{i=0}^m\pi_i^{-1}[\{0\}],
\]
so $\varphi$ is a homeomorphic isomorphism of groups, $\hat G\cong\varprojlim G/G_n$. From now on we will denote by $\hat G$ either $\hat G$ or $\varprojlim G/G_n$ and we will use both the analytical and algebraic description. 

For example, it is easy to see that $\hat G_n=\ker(\hat G\to G/G_n)$.

\subsection{Completeness}

Recall that

\begin{defi}
    Given a topological abelian group $G$, we say that it is complete if all Cauchy sequences converge.
\end{defi}

Let $(g_n)_{n\geq 0}$ be a convergent sequence to $g\in G$; i.e., for all neighborhood of $0$, $U$, there exists an $n_U\geq 0$ such that for all $n\geq n_U$ we have $g_n\in g-U$. That is, $g-g_n\in U$. Since $G$ is a topological group we have that $(g_n)_{n\geq 0}\in G^\mathbb N_\mathcal C$. So in a complete topological abelian group all Cauchy sequences converge and all convergent sequences are Cauchy.

We end this section giving an alternative definition of completeness of Hausdorff topological abelian groups and proving that the Hausdorff abelian group $\hat G$ is complete.
\begin{thm}\label{completeness}
    Given a topological abelian group $G$. The homomorphism $\phi\colon G\to \hat G$ is an isomorphism if and only if $G$ is Hausdorff and complete.
\end{thm}
\begin{proof}
    We already know that $\phi$ is one-to-one iff $G$ is Hausdorff. On the other hand, $\phi$ is onto iff every Cauchy sequence is equivalent to a constant sequence iff every sequence converges iff $G$ is complete.
\end{proof}

From now on the algebraic description of $\hat G$ will come in handy to prove its completeness.

Let $(A_n)_{n\geq 0}$ be an inverse system and $A:=\prod_{n\geq 0}A_n$. The map
\[
\begin{matrix}
    d^A\colon &A&\longrightarrow& A\\
    &(a_n)_{n\geq 0} &\longmapsto & (a_n-\theta_{n+1}(a_{n+1}))_{n\geq 0}
\end{matrix}
\]
is a homomorphism of abelian groups such that $\ker d^A=\varprojlim A_n$. If each $\theta_n$ is surjective; i.e., $(A_n)_{n\geq 0}$ is surjective, then $d^A$ is also surjective. 

\begin{defi}
    Given $(A_n)_{n\geq 0}$, $(B_n)_{n\geq 0}$ and $(C_n)_{n\geq 0}$ three inverse systems. If for each $n\geq 0$ we have the following commutative diagram of exact sequences
    \[
    \begin{tikzcd}
        0 \ar[r] & A_{n+1} \ar[r]\ar[d] & B_{n+1} \ar[r]\ar[d] & C_{n+1} \ar[r]\ar[d] & 0\\
        0 \ar[r] & A_n \ar[r] & B_n \ar[r] & C_n \ar[r] & 0,
    \end{tikzcd}
    \]
    we say that we have an \emph{exact sequence of inverse systems} denoted by
    \[
    \begin{tikzcd}
        0 \ar[r] & (A_n)_{n\geq 0} \ar[r] & (B_n)_{n\geq 0} \ar[r] & (C_n)_{n\geq 0} \ar[r] &0.
    \end{tikzcd}
    \]
\end{defi}

\begin{prop}\label{exactseq}
    Let $0\to (A_n)_{n\geq 0}\to (B_n)_{n\geq 0}\to(C_n)_{n\geq 0}\to 0$ be an exact sequence of inverse systems then
    \[
    \begin{tikzcd}
        0 \ar[r] & \varprojlim A_n \ar[r] &  \varprojlim B_n \ar[r] & \varprojlim C_n 
    \end{tikzcd}
    \]
    is exact. Also, if $(A_n)_{n\geq 0}$ is a surjective system,
    \[
    \begin{tikzcd}
        0 \ar[r] & \varprojlim A_n \ar[r] &  \varprojlim B_n \ar[r] & \varprojlim C_n \ar[r] & 0 
    \end{tikzcd}
    \]
    is exact.
\end{prop}
\begin{proof}
    We have the following commutative diagram of exact sequences
    \[
    \begin{tikzcd}
         0 \ar[r] & \prod_{n\geq 0}A_{n} \ar[r]\ar[d, "d^A"] & \prod_{n\geq 0} B_{n} \ar[r]\ar[d, "d^B"] & \prod_{n\geq 0}C_{n} \ar[r]\ar[d, "d^C"] & 0\\
        0 \ar[r] & \prod_{n\geq 0}A_n \ar[r] & \prod_{n\geq 0}B_n \ar[r] & \prod_{n\geq 0}C_n \ar[r] & 0.
    \end{tikzcd}
    \]
    Thus, by the snake lemma, we get an exact sequence
    \[
    \begin{tikzcd}
       0 \rar & \ker d^A \rar & \ker d^B \rar
             \ar[draw=none]{d}[name=X, anchor=center]{}
    & \ker d^C \ar[rounded corners,
            to path={ -- ([xshift=2ex]\tikztostart.east)
                      |- (X.center) \tikztonodes
                      -| ([xshift=-2ex]\tikztotarget.west)
                      -- (\tikztotarget)}]{dll}[at end]{} \\      
  &\coker d^A \rar & \coker d^B \rar & \coker d^C \rar & 0.
    \end{tikzcd}
    \]
\end{proof}

\begin{coro}\label{corohat}
    Let $G$ be a topological abelian group with topology defined by a sequence $(G_n)_{n\geq 0}$. If $0\to H \overset{\iota}{\to} G\overset{\pi}{\to} K\to 0$ is exact then 
    \[
    \begin{tikzcd}
        0 \rar & \hat H \rar & \hat G \rar & \hat K \rar & 0
    \end{tikzcd}
    \]
    is exact, where the topologies of $H$ and $K$ are given by the sequences $(\iota^{-1}[G_n])_{n\geq 0}$ and $(\pi[G_n])_{n\geq 0}$ respectively.
\end{coro}
\begin{proof}
    It is clear that $(H/\iota^{-1}[G_n])_{n\geq 0}$ is a surjective inverse system and 
    \[
    \begin{tikzcd}
        0 \rar & (H/\iota^{-1}[G_n])_{n\geq 0} \rar & (G/G_n)_{n\geq 0}\rar & (K/\pi[G_n])_{n\geq 0}\rar & 0
    \end{tikzcd}
    \]
    is exact. Now apply proposition \ref{exactseq}.
\end{proof}

Recall that $G/G_n$ has the discrete topology as a topological abelian subgroup of $G$. Hence $G/G_n$ is Hausdorff ---so $\phi\colon G/G_n \hookrightarrow \widehat{G/G_n}$--- and every Cauchy sequence is ultimately constant in $G/G_n$ so that $\phi$ is onto. Thus $\widehat{G/G_n}\cong G/G_n$. Applying corollary \ref{corohat} for $H=G_n$ and $K=G/G_n$ we get
\[
\hat G/\hat G_n\cong\widehat{G/G_n}\cong G/G_n.
\]
Taking inverse limits this becomes
\begin{coro}
    $\hat{\hat{G}}\cong \hat G$.
\end{coro}

So $\hat G$ is Hausdorff and complete: $\phi\colon \hat G\to \hat{\hat {G}}$ is one-to-one and onto and apply theorem \ref{completeness}.


